\documentclass[10pt]{article}
\usepackage[margin=1in]{geometry}
\usepackage{enumitem}
\usepackage{hyperref}
\usepackage{xcolor}
\hypersetup{colorlinks=true, urlcolor=blue, linkcolor=black}
\usepackage{titlesec}
\titleformat{\section}{\normalfont\large\bfseries}{\thesection}{0.6em}{}
\titlespacing*{\section}{0pt}{1.2ex}{0.6ex}
\setlist{nosep, leftmargin=1.4em}

\newcommand{\code}[1]{\texttt{#1}}

\title{\vspace{-2.5em}\textbf{Artifact Appendix:\\ Interval POMDP Shielding for Imperfect-Perception Agents}}
\author{William Scarbro \and Ravi Mangal}
\date{EMSOFT 2026 --- Artifact Evaluation}

\begin{document}
\maketitle
\vspace{-1.5em}

\section*{Abstract}
This artifact reproduces every figure and table in \emph{Interval POMDP Shielding
for Imperfect-Perception Agents}. It contains the experiment code, the trained
DQN controllers, the optimized adversarial perception realizations, and the
bundled result data, together with a Docker recipe for a self-contained,
CPU-only environment. The artifact evaluates five runtime shields
(\textbf{Envelope}, \textbf{Single-Belief}, \textbf{Observation}, \textbf{Carr},
\textbf{Fwd-Sampling}) on four benchmarks (\textbf{TaxiNet}, \textbf{Obstacle},
\textbf{CartPole}, \textbf{Refuel}). Everything runs through two commands:
a fast \code{smoke\_test} (${\sim}2$\,min) that regenerates all figures/tables
from bundled data, and a full \code{reproduce\_results} (${\sim}12$\,h) that
re-executes the underlying Monte Carlo experiments.

\section*{Artifact check-list (meta-information)}
\begin{itemize}
  \item \textbf{Algorithms:} interval-POMDP runtime shielding; LP envelope
        over-approximation; forward sampling; Cross-Entropy Method for
        adversarial realization optimization; Clopper--Pearson confidence intervals.
  \item \textbf{Program:} Python 3.11 package \code{ipomdp\_shielding}.
  \item \textbf{Model:} bundled DQN controllers (\code{results/cache/*\_agent.pt}); no retraining required.
  \item \textbf{Data set:} bundled result JSONs, perception realizations, and
        TaxiNet point-estimate CSVs (all included; no external download).
  \item \textbf{Run-time environment:} Docker (base image \code{python:3.11-slim}); CPU-only; no network at run time.
  \item \textbf{Hardware:} any x86-64 machine; ${\ge}8$\,GB RAM; ${\sim}1$\,GB disk. No GPU.
  \item \textbf{Metrics:} failure / stuck / safe rates; per-step shield latency; envelope coarseness.
  \item \textbf{Output:} figures written to \code{figures/}; tables printed to stdout.
  \item \textbf{Experiments:} \code{./run\_experiments.sh smoke\_test} and \code{reproduce\_results}.
  \item \textbf{How much disk space required?} ${\sim}1$\,GB (image + artifact).
  \item \textbf{How much time to prepare workflow?} ${\sim}5$\,min (\code{docker build}).
  \item \textbf{How much time to complete experiments?} \code{smoke\_test} ${\sim}2$\,min; \code{reproduce\_results} ${\sim}12$\,h.
  \item \textbf{Publicly available?} Yes. \item \textbf{Code license:} MIT.
  \item \textbf{Archived (DOI):} \href{https://doi.org/10.5281/zenodo.21539082}{10.5281/zenodo.21539082}.
\end{itemize}

\section*{Description}
\subsection*{How to access}
Archived on Zenodo at \href{https://doi.org/10.5281/zenodo.21539082}{10.5281/zenodo.21539082}
(MIT-licensed). Source is also at
\url{https://github.com/WilliamScarbro/ipomdp_shielding_artifact}. The accepted
paper is bundled with the artifact as \code{paper.pdf}.

\subsection*{Technical skills required to review the artifact}
Reviewing the artifact requires only basic command-line and Docker proficiency:
building an image (\code{docker build}) and running a container with a volume
mount (\code{docker run -v}). No GPU, cluster, licensed software, or special
hardware is needed, and the artifact requires no network access at run time.
Interpreting the regenerated figures and tables benefits from --- but does not
require --- familiarity with POMDPs, interval/robust Markov models, and runtime
shielding; the mapping from each command to its paper figure/table is given in
\code{README.md} so that outputs can be checked against the paper without domain
expertise.

\subsection*{Hardware dependencies}
Any x86-64 machine. CPU-only --- no GPU. ${\ge}8$\,GB RAM (the figure path uses
well under 2\,GB) and ${\sim}1$\,GB free disk.

\subsection*{Software dependencies}
Docker Engine only; the image pins the full environment. For reference, the
reproduction path depends solely on \code{numpy 1.26.4}, \code{scipy 1.13.1},
\code{statsmodels 0.14.2}, \code{matplotlib 3.9.2}, and \code{torch 2.4.1} (CPU
build), as pinned in \code{requirements.txt}.

\section*{Installation}
From the artifact root:
\begin{verbatim}
docker build -t ipomdp-shielding .
\end{verbatim}

\section*{Experiment workflow}
Both commands write the paper figures into \code{figures/}; mount a host directory there.
\begin{verbatim}
# Fast check + regenerate all figures/tables from bundled data (default CMD, ~2 min):
docker run --rm -v "$PWD/figures:/artifact/figures" ipomdp-shielding

# Full reproduction: rerun all Monte Carlo experiments, then replot (~12 h):
docker run --rm -v "$PWD/figures:/artifact/figures" ipomdp-shielding \
    ./run_experiments.sh reproduce_results
\end{verbatim}
\code{smoke\_test} finishes with \code{[smoke] ALL CHECKS PASSED}. \code{reproduce\_results}
reuses the bundled trained agents and optimized realizations (no retraining) and
overwrites the result JSONs under \code{results/} before replotting.

\section*{Evaluation and expected results}
After either command, \code{figures/} contains the regenerated plots and the
supplement tables are printed to stdout:
\begin{itemize}
  \item Fig.~1/2 --- main summary bars: \code{summary\_v7\_bars.png}, \code{summary\_v7\_safe\_bars.png}
  \item Fig.~3 --- Pareto: \code{pareto\_v7\_taxinet.png}, \code{pareto\_v7\_obstacle.png}
  \item Fig.~4 --- conformal comparison: \code{taxinet\_v2\_extremal\_comparison.png} (+ Table~3 to stdout)
  \item Shield timing: \code{inference\_timing\_summary.pdf}
  \item Coarseness: \code{coarse\_taxinet\_results.png}; Perception variability: \code{perception\_variability\_taxinet.png}
  \item Alpha sweep: \code{alpha\_sweep\_taxinet\_\{fail,stuck,safe\}.png}
  \item Tables~1--3 printed to stdout.
\end{itemize}
Regenerated figures and printed tables match the paper; e.g.\ the conformal
comparison (Fig.~4 / Table~3) reproduces exactly and the summary/Pareto figures
are byte-for-byte equivalent to the submission. A full command-to-figure mapping
is in \code{README.md}.

\section*{Badges claimed}
\textbf{Available} (public Zenodo DOI, MIT license), \textbf{Reviewed/Functional}
(\code{smoke\_test} regenerates all outputs and exercises the live shielding
pipeline), and \textbf{Reproducible} (\code{reproduce\_results} recomputes the
underlying results). See \code{STATUS.md} for per-badge justification.

\end{document}
